% FAC-EN-U0079
% Source: fac_body.tex lines 3358--3400.
% Authority: PDF physical pp. 81--82; printed pp. 276--277
% (no. 81 complete; hard stop before Paris and the bibliography).

\medskip
\phantomsection\label{fac:III:6:no:81}
\textbf{81. Degree of the characteristic function.}
If $\mathfrak F$ is a coherent algebraic sheaf on an algebraic variety
$V$, we call the set of points $x\in V$ such that
$\mathfrak F_x\neq0$ the \phantomsection\label{fac:III:6:no:81:def:support}
\textit{support} of $\mathfrak F$, and denote it by
$\operatorname{Supp}(\mathfrak F)$. Since $\mathfrak F$ is a sheaf of
finite type, this set is \textit{closed}: indeed, if
$\mathfrak F_x=0$, the zero section generates $\mathfrak F_x$, and
hence also $\mathfrak F_y$ for every $y$ sufficiently close to $x$ by
\hyperref[fac:I:2:no:12:prop:1]{Proposition~1 of no.~12}. Thus the
complement of $\operatorname{Supp}(\mathfrak F)$ is open.

Let $M$ be a finitely generated graded $S$-module, and let
$\mathfrak F=\mathfrak A(M)$ be the sheaf defined by $M$ on
$\mathbf P_r(K)=X$. We can determine $\operatorname{Supp}(\mathfrak F)$
from $M$ as follows.

Let
\phantomsection\label{fac:III:6:no:81:primary-decomposition}
\[
0=\bigcap_\alpha M^\alpha
\]
be a decomposition of $0$ as an intersection of homogeneous primary
submodules $M^\alpha$ of $M$, with the $M^\alpha$ corresponding to
homogeneous prime ideals $\mathfrak p^\alpha$ (cf.~[12], Chap.~IV).
We shall suppose that this decomposition is ``as short as possible,''
that is, that no $M^\alpha$ is contained in the intersection of the
others. For every $x\in X$, each $\mathfrak p^\alpha$ defines a prime
ideal $\mathfrak p_x^\alpha$ of the local ring $\mathcal O_x$, and
$\mathfrak p_x^\alpha=\mathcal O_x$ if and only if $x$ does not belong
to the variety $V^\alpha$ defined by $\mathfrak p^\alpha$. Likewise,
\phantomsection\label{fac:III:6:no:81:localized-primary-decomposition}
\[
0=\bigcap_\alpha M_x^\alpha
\]
in $M_x$, and one readily verifies that this gives a primary
decomposition of $0$ in $M_x$, with the $M_x^\alpha$ corresponding to
the prime ideals $\mathfrak p_x^\alpha$. If $x\notin V^\alpha$, then
$M_x^\alpha=M_x$; restricting to the $M_x^\alpha$ for which
$x\in V^\alpha$ gives a decomposition ``as short as possible''
(cf.~[12], Chap.~IV, th.~4, where analogous results are established).
It follows immediately that $M_x\neq0$ if and only if $x$ belongs to
one of the varieties $V^\alpha$; in other words,
\phantomsection\label{fac:III:6:no:81:formula:support-union}
\[
\operatorname{Supp}(\mathfrak F)=\bigcup_\alpha V^\alpha.
\]

\phantomsection\label{fac:III:6:no:81:prop:6}
\textsc{Proposition 6.} \textit{If $\mathfrak F$ is a coherent
algebraic sheaf on $\mathbf P_r(K)$, the degree of the polynomial
$\chi(\mathfrak F(n))$ is equal to the dimension of
$\operatorname{Supp}(\mathfrak F)$.}

\phantomsection\label{fac:III:6:no:81:prop:6:proof}
We argue by induction on $r$, the case $r=0$ being trivial. We may
suppose that $\mathfrak F=\mathfrak A(M)$, where $M$ is a finitely
generated graded $S$-module. With the
\hyperref[fac:III:6:no:81:primary-decomposition]{notation introduced above}, we
must prove that $\chi(\mathfrak F(n))$ is a polynomial of degree
\phantomsection\label{fac:III:6:no:81:formula:q}
\[
q=\sup_\alpha\dim V^\alpha.
\]

Let $t$ be a homogeneous linear form belonging to none of the prime
ideals $\mathfrak p^\alpha$, except possibly the ``improper'' prime
ideal $\mathfrak p^0=(t_0,\ldots,t_r)$; such a form exists because the
field $K$ is infinite. Let $E$ be the hyperplane of $X$ with equation
$t=0$. Consider the exact sequence
\phantomsection\label{fac:III:6:no:81:sequence:hyperplane}
\[
0\longrightarrow\mathcal O(-1)\longrightarrow\mathcal O
\longrightarrow\mathcal O_E\longrightarrow0,
\]
where $\mathcal O\to\mathcal O_E$ is the restriction homomorphism,
while $\mathcal O(-1)\to\mathcal O$ is the homomorphism
$f\mapsto t\cdot f$. Tensoring with $\mathfrak F$ gives the exact
sequence
\phantomsection\label{fac:III:6:no:81:sequence:tensored}
\[
\mathfrak F(-1)\longrightarrow\mathfrak F
\longrightarrow\mathfrak F_E\longrightarrow0,
\qquad
\mathfrak F_E=\mathfrak F\otimes_{\mathcal O}\mathcal O_E.
\]
Over $U_i$, we may identify $\mathfrak F(-1)$ with $\mathfrak F$, and
under this identification the homomorphism
$\mathfrak F(-1)\to\mathfrak F$ just defined becomes multiplication
by $t/t_i$. Since $t$ was chosen outside the $\mathfrak p^\alpha$,
$t/t_i$ belongs to none of the prime ideals of
$M_x=\mathfrak F_x$ when $x\in U_i$, and the
\hyperref[fac:III:6:no:81:sequence:tensored]{preceding homomorphism} is
injective (cf.~[12], p.~122, th.~7, b$'''$). We therefore have the
exact sequence
\phantomsection\label{fac:III:6:no:81:sequence:injective}
\[
0\longrightarrow\mathfrak F(-1)\longrightarrow\mathfrak F
\longrightarrow\mathfrak F_E\longrightarrow0,
\]
and hence, for every $n\in\mathbf Z$, the exact sequence
\phantomsection\label{fac:III:6:no:81:sequence:twisted}
\[
0\longrightarrow\mathfrak F(n-1)\longrightarrow\mathfrak F(n)
\longrightarrow\mathfrak F_E(n)\longrightarrow0.
\]
Applying \hyperref[fac:III:6:no:79:prop:1]{Proposition~1}, we obtain
\phantomsection\label{fac:III:6:no:81:formula:first-difference}
\[
\chi(\mathfrak F(n))-\chi(\mathfrak F(n-1))
=\chi(\mathfrak F_E(n)).
\]

But $\mathfrak F_E$ is a coherent sheaf of $\mathcal O_E$-modules;
in other words, it is a coherent algebraic sheaf on $E$, which is a
projective space of dimension $r-1$. Moreover,
$\mathfrak F_{x,E}=0$ means that the endomorphism of $\mathfrak F_x$
defined by multiplication by $t/t_i$ is surjective, which implies
$\mathfrak F_x=0$ (cf.~[6], Chap.~VIII, prop.~5.1$'$). Consequently,
\phantomsection\label{fac:III:6:no:81:formula:support-intersection}
\[
\operatorname{Supp}(\mathfrak F_E)
=E\cap\operatorname{Supp}(\mathfrak F).
\]
Since $E$ contains none of the varieties $V^\alpha$, it follows from
a known result that $\dim\operatorname{Supp}(\mathfrak F_E)=q-1$. By
the induction hypothesis, $\chi(\mathfrak F_E(n))$ is therefore a
polynomial of degree $q-1$; since it is the first difference of the
function $\chi(\mathfrak F(n))$, the latter is indeed a polynomial of
degree $q$.

\textsc{Remarks}. \phantomsection\label{fac:III:6:no:81:remark:1}
(1) \hyperref[fac:III:6:no:81:prop:6]{Proposition~6} was well known when
$\mathfrak F=\mathcal O/\mathcal J$, with $\mathcal J$ a coherent
sheaf of ideals. See, for example, [9], no.~24.

\phantomsection\label{fac:III:6:no:81:remark:2}
(2) The \hyperref[fac:III:6:no:81:prop:6:proof]{preceding proof} does not use
\hyperref[fac:III:6:no:80:prop:3]{Proposition~3} and therefore
proves it anew.
